\documentclass[11pt]{article}
\usepackage[a4paper,margin=27mm]{geometry}
\usepackage[T1]{fontenc}
\usepackage[utf8]{inputenc}
\usepackage{lmodern,microtype}
\usepackage{amsmath,amssymb,amsthm,mathtools,bm}
\usepackage{booktabs,array,enumitem}
\usepackage{xcolor}
\usepackage[colorlinks=true,linkcolor=blue!45!black,citecolor=blue!45!black,urlcolor=blue!55!black]{hyperref}
\usepackage[nameinlink,noabbrev]{cleveref}
\usepackage[backend=biber,style=numeric,sorting=nyt,maxbibnames=99]{biblatex}
\addbibresource{references.bib}
\AtBeginBibliography{\sloppy}
\setlength{\emergencystretch}{2em}

\newtheorem{theorem}{Theorem}[section]
\newtheorem{lemma}[theorem]{Lemma}
\newtheorem{proposition}[theorem]{Proposition}
\newtheorem{corollary}[theorem]{Corollary}
\newtheorem{remark}[theorem]{Remark}
\newcommand{\E}{\mathbb E}
\newcommand{\R}{\mathbb R}
\newcommand{\Sym}{\operatorname{Sym}}
\newcommand{\Herm}{\operatorname{Herm}}
\newcommand{\tr}{\operatorname{tr}}
\newcommand{\dist}{\operatorname{dist}}
\newcommand{\supp}{\operatorname{supp}}
\newcommand{\cC}{\mathcal C}
\newcommand{\cT}{\mathcal T}
\newcommand{\ip}[2]{\langle #1,#2\rangle}
\newcommand{\norm}[1]{\lVert #1\rVert}

\hypersetup{
 pdftitle={Finite-Sample-Valid Tests for Structured Matricial Hausdorff Moments: Joint Gaussian Grams, Strict Half-Time Separation, and Sharp Tangent-Cone Power},
 pdfauthor={Lluis Eriksson},
 pdfkeywords={matrix moment problem, Hausdorff moments, Wishart, Loewner confidence set, tangent cone, exact certificate}
}

\title{\bfseries Finite-Sample-Valid Tests for Structured Matricial Hausdorff Moments\\[3pt]
\large Joint Gaussian Grams, Strict Half-Time Separation,\\
and Sharp Tangent-Cone Power}
\author{Lluis Eriksson}
\date{August 14, 2026}

\begin{document}
\maketitle

\begin{abstract}
We give a finite-sample test for whether one joint Gaussian sketch covariance
is compatible with a positive-semidefinite matrix-valued measure supported on
a prescribed interval.  The observed vector contains several exponent
blocks, so its population covariance is one block-Hankel Gram matrix rather
than a list of unrelated moment estimates.  The classical truncated
matricial Hausdorff theorem supplies the exact parity-dependent null.
A single Gaussian singular-value event gives a simultaneous Loewner band for
the whole covariance; intersecting that band with the structured moment cone
yields a nonasymptotic level-$\alpha$ semidefinite test without sample
splitting.

The full joint Gram is strictly more informative than the integer-time
localizer used in the preceding finite-sample method.  For
$\nu=(1-w)\delta_0+w\delta_t$, with $0<s<t\le1$ and
$w<s^2/t^2$, the previous condition $s^2M_0-M_2>0$ holds,
whereas the half-time condition $sM_1-M_2=wt(s-t)<0$.  We prove that
the old test has power at most $\alpha$, while the full test has power at
least $1-\alpha$ above an explicit sample threshold on the same acquisition.
At a positive-definite covariance on a singular null face, the Wishart
experiment is locally asymptotically normal.  The constrained likelihood
ratio converges to squared Gaussian distance from an explicit spectrahedral
tangent cone, giving its pointwise local power function and a matching
$n^{-1/2}$ separation boundary.  We give auditable dual semantics and an
exact rational certificate for the strict fixture.

The moment characterization, Gaussian covariance bounds, and generic
constrained-likelihood theory are established ingredients.  The contribution
is their structured joint-sketch integration, the analytic same-data
separation, and the exact certificate layer.  This work explicitly extends
ARR record \texttt{ARR-2026-77QM18J2KG9679B7}.
\end{abstract}

\section{Introduction}

Finite moment tests turn support claims into positive-semidefinite
constraints, but several shortcuts destroy their meaning.  A single scalar
localizer does not test the complete truncated moment cone; independently
estimated moments ignore the cross-covariances in a joint sketch; and a
negative floating-point eigenvalue is not a certificate.  These distinctions
are most important on singular faces, precisely where support restrictions
become informative.

The matrix Hausdorff moment problem on a compact interval is classical.
Parity-complete block-Hankel characterizations and matrix canonical moments
are developed in
\cite{DetteStudden2002,ChoqueRiveroEtAl2006,ChoqueRiveroEtAl2007,FritzscheKirsteinMaedler2019}.
Moment and localizing matrices are standard in polynomial optimization
\cite{Lasserre2001,HeltonNie2012}.  We specialize those results without
discarding the off-diagonal moment information present in the experiment.

The statistical starting point is the previous record \cite{ARRPrior2026}.
Its integer-time support component retains
\[
 [\theta B_{i+j}-B_{i+j+1}]_{i,j=0}^N\succeq0.                    \tag{1.1}
\]
Here we put $y=\sqrt{x}$ and retain the half-time covariance blocks already
acquired.  The whole covariance becomes a block-Hankel matrix of $y$-moments,
and the correct support localizer uses every polynomial direction on that
grid.  No extra sample or multiple-testing correction is needed: all
constraints are optimized inside one simultaneous band.

Our contributions are: an operational joint Gaussian model; a parity-exact
structured test; an analytic old-versus-full power separation; consistency
against every fixed covariance alternative; an active-kernel tangent formula
and Gaussian cone-projection likelihood limit; and exact dual semantics.
Rejection falsifies compatibility of the visible matrix measure with the
declared support.  Non-rejection does not prove a physical spectral gap.

\section{Joint Gaussian exponent sketches}

Let $\mathsf T$ be a positive contraction on a real Hilbert space and let
$\psi_1,\ldots,\psi_c$ be fixed probes.  For $k=0,\ldots,q$, assume
$\mathsf T^{k/2}\psi_a$ is defined.  Let
$W=\{W(h):h\in\mathcal H\}$ be an isonormal Gaussian process, so
$\E W(h)W(k)=\ip{h}{k}$, and set
\[
 z_{k,a}=W(\mathsf T^{k/2}\psi_a),\qquad
 z=(z_{0,1},\ldots,z_{q,c})\in\R^d,\quad d=(q+1)c.                 \tag{2.1}
\]
In finite dimension one may realize $W(h)=\ip{g}{h}$ with
$g\sim N(0,I)$.  If $E_{\mathsf T}$ is the spectral resolution of
$\mathsf T$, define the real-symmetric matrix measure
\[
 \bm\nu_{ab}(B)=\ip{\psi_a}{E_{\mathsf T}(\{\lambda:\sqrt\lambda\in B\})\psi_b}.
\]
The vector is jointly Gaussian of known mean zero, and its covariance blocks
are
\[
 \E z_{k,a}z_{\ell,b}
 =\ip{\psi_a}{\mathsf T^{(k+\ell)/2}\psi_b}
 =:(M_{k+\ell})_{ab}.                                             \tag{2.2}
\]
Consequently
\[
 K:=\E zz^\mathsf T=H_q(M):=[M_{i+j}]_{i,j=0}^q.                  \tag{2.3}
\]
The spectral theorem gives a positive-semidefinite $c\times c$
matrix-valued measure $\bm\nu$ on $[0,1]$ with
\[
 M_j=\int_0^1 y^j\,d\bm\nu(y).                                    \tag{2.4}
\]
The visible support claim is $\supp\bm\nu\subset[0,s]$.

Observe iid copies $z_1,\ldots,z_n$ and form
\[
 S=\frac1n\sum_{r=1}^n z_rz_r^\mathsf T.                           \tag{2.5}
\]
Sampling noise makes the empirical anti-diagonal blocks unequal.  We do not
project $S$ onto the Hankel subspace.  Instead, the structured population
moments remain primal variables and $H_q(M)$ is required to lie in a
confidence band around $S$.  This preserves coverage independently of any
projection rule.

\begin{remark}
The blocks $M_j$ are not independently sampled covariances: they are repeated
anti-diagonals of one $K$.  The off-diagonal blocks contain the half-time
moments.  Replacing them by independent moment oracles destroys the same-data
comparison proved below.
\end{remark}

\section{The parity-complete matricial cone}

For real-symmetric blocks define
\[
\begin{aligned}
 H_r(M)&=[M_{i+j}]_{i,j=0}^r,&
 G_r(M)&=[M_{i+j+1}]_{i,j=0}^r,\\
 L_{r,s}(M)&=[sM_{i+j}-M_{i+j+1}]_{i,j=0}^r,&
 C_{r,s}(M)&=[sM_{i+j+1}-M_{i+j+2}]_{i,j=0}^r.
\end{aligned}                                                       \tag{3.1}
\]

\begin{theorem}[truncated matricial Hausdorff criterion]\label{thm:cone}
The blocks $M_0,\ldots,M_{2r}$ have a positive matrix-valued representing
measure on $[0,s]$ if and only if
\[
 H_r(M)\succeq0,\qquad C_{r-1,s}(M)\succeq0,                      \tag{3.2}
\]
with the second condition absent at $r=0$.  The blocks through $M_{2r+1}$
have such a measure if and only if
\[
 G_r(M)\succeq0,\qquad L_{r,s}(M)\succeq0.                        \tag{3.3}
\]
Both statements include rank-deficient sequences.
\end{theorem}

\begin{proof}
Necessity follows from vector-polynomial Gram identities.  For
$p(y)=\sum_{i=0}^ry^iv_i$,
\[
 v^*H_r(M)v=\int p(y)^*\,d\bm\nu(y)\,p(y)\ge0,                    \tag{3.4}
\]
and the remaining matrices insert $y$, $s-y$, or $y(s-y)$.

For sufficiency, define
$\mathcal L_M(\sum_jF_jy^j)=\sum_j\tr(F_jM_j)$.  Every real-symmetric matrix
polynomial nonnegative on $[0,s]$ of degree at most $2r$ factors as
\[
 F=A^\mathsf TA+y(s-y)B^\mathsf TB,\qquad \deg A\le r,\quad\deg B\le r-1,          \tag{3.5}\label{eq:evenfactor}
\]
and at degree at most $2r+1$ as
\[
 F=yA^\mathsf TA+(s-y)B^\mathsf TB,\qquad \deg A,\deg B\le r.                     \tag{3.6}\label{eq:oddfactor}
\]
Here $A$ and $B$ may be rectangular matrix polynomials; equivalently, each
term denotes a finite sum of real-symmetric polynomial squares.
Thus the displayed LMIs make $\mathcal L_M$ nonnegative on every positive
matrix polynomial of the declared degree.

The truncated matrix-measure cone is closed.  Bounded $M_0$ bounds total
scalar mass; positivity bounds the variation of every entry; weak-star
compactness supplies a matrix-measure limit.  If $M$ lay outside that cone,
strict separation would produce a positive matrix polynomial with negative
$\mathcal L_M$ value, contradicting \cref{eq:evenfactor,eq:oddfactor}.
\end{proof}

The theorem is established literature, not our novelty claim.  Entrywise
scalar tests are insufficient because independently chosen scalar measures
need not assemble into one positive matrix measure.

For the acquisition $q=2N+1$, the covariance exposes
$M_0,\ldots,M_{2q}=M_{4N+2}$.  The exact even null is
\[
 K=H_q(M),\qquad C_{q-1,s}(M)\succeq0.                             \tag{3.7}
\]
The Loewner band already enforces $K\succeq0$; the new ingredients are the
anti-diagonal identities and full half-time localizer.

\section{A structured Loewner confidence test}

Let $G$ be a $d\times n$ standard real Gaussian matrix.  On the range of
$K$, the sample covariance has the congruence representation
\[
 S=K^{1/2}\left(\frac1nGG^\mathsf T\right)K^{1/2}.                 \tag{4.1}
\]
For $0<\alpha<1$, put
\[
 \eta=\frac{\sqrt d+\sqrt{2\log(2/\alpha)}}{\sqrt n},              \tag{4.2}
\]
assume $\eta<1$, and define
\[
 a=(1+\eta)^{-2},\qquad b=(1-\eta)^{-2}.                           \tag{4.3}
\]
The Gaussian singular-value inequality
\cite{DavidsonSzarek2001,KoltchinskiiLounici2017} yields
\[
 \Pr_K\{aS\preceq K\preceq bS\}\ge1-\alpha.                       \tag{4.4}\label{eq:coverage}
\]
This is a nonasymptotic valid Gaussian Loewner band.  It is generally
conservative, so we do not call it an exact-coverage Wishart interval.
Exact Wishart eigenvalue quantiles may replace $(a,b)$ without changing the
primal geometry.

Given $S$, solve
\[
\begin{array}{ll}
\text{find}&M_0,\ldots,M_{2q}\in\Sym_c\\[1mm]
\text{such that}&aS\preceq H_q(M)\preceq bS,\\
&C_{q-1,s}(M)\succeq0.
\end{array}                                                        \tag{4.5}
\]
Reject exactly when this structured feasibility problem is infeasible.
Using moment blocks as variables enforces all anti-diagonal equalities by
construction and makes the adjoint map auditable.

\begin{theorem}[finite-sample validity]\label{thm:level}
Under the joint Gaussian model,
\[
 \sup_{\supp\bm\nu\subset[0,s]}
 \Pr_{\bm\nu}\{\text{reject}\}\le\alpha.                           \tag{4.6}
\]
The optimization may inspect every parity-appropriate constraint and matrix
witness formed from the acquired truncation $M_0,\ldots,M_{2q}$ adaptively.
No further multiplicity correction is required.
\end{theorem}

\begin{proof}
On the single event \cref{eq:coverage}, the true blocks satisfy the band with
$H_q(M)=K$.  Under the null, \cref{thm:cone} gives
$C_{q-1,s}(M)\succeq0$.  The true sequence is feasible, so rejection is
contained in the complement of the coverage event.
\end{proof}

\begin{remark}[singular covariance]
The level argument does not invert $K$ and is valid on singular covariance
faces.  Implementations should compress to the exact common range implied by
the band or use facial reduction.  Adding a numerical ridge changes the null.
\end{remark}

\section{Strict same-data power separation}\label{sec:strict}

Consider the scalar measure
\[
 \nu=(1-w)\delta_0+w\delta_t,\qquad
 0<s<t\le1,\quad0<w<s^2/t^2.                                     \tag{5.1}\label{eq:twoatom}
\]
Its first moments and joint $N=0$ covariance are
\[
 M_0=1,\quad M_1=wt,\quad M_2=wt^2,\qquad
 K=\begin{pmatrix}1&wt\\wt&wt^2\end{pmatrix}\succ0.               \tag{5.2}
\]
The preceding integer-time method uses $x=y^2$, $\theta=s^2$, and retains
\[
 \theta M_0-M_2=s^2-wt^2>0.                                      \tag{5.3}\label{eq:oldmargin}
\]
Its finite-window population null is therefore true despite $t>s$.  The
complete half-time null contains instead
\[
 C_{0,s}=sM_1-M_2=wt(s-t)<0.                                     \tag{5.4}
\]
Both tests use the same $K$; the predecessor discards its off-diagonal
$M_1$.

For a self-contained comparison, define the preceding $N=0$ test on the
same observed $S$ as follows: accept whenever there is a symmetric
$Q=\begin{psmallmatrix}B_0&u\\u&B_1\end{psmallmatrix}$ satisfying
$aS\preceq Q\preceq bS$ and $s^2B_0-B_1\ge0$.  Thus it retains the two
integer-time diagonal blocks and leaves the half-time cross block $u$
unconstrained.  On the event \cref{eq:coverage}, the true $Q=K$ is feasible
for this old test whenever \cref{eq:oldmargin} holds.

\begin{theorem}[strict old-versus-full separation]\label{thm:strict}
For \cref{eq:twoatom}, the preceding $N=0$ integer-time test has rejection
probability at most $\alpha$ whenever its level theorem applies.  Define
\[
 \gamma=wt(t-s),\qquad
 A=\begin{pmatrix}0&s/2\\s/2&-1\end{pmatrix},                     \tag{5.5}
\]
\[
 W=\tr|K^{1/2}AK^{1/2}|
 =\sqrt{\gamma^2+s^2w(1-w)t^2},\qquad r=\gamma/W.                 \tag{5.6}
\]
The full structured test has rejection probability at least $1-\alpha$
whenever
\[
 \boxed{
 n>
 \bigl(\sqrt2+\sqrt{2\log(2/\alpha)}\bigr)^2
 \frac{(\sqrt{1+r}+1)^4}{r^2}.}                                   \tag{5.7}\label{eq:threshold}
\]
The scalar construction embeds algebraically in every channel dimension by
tensoring with a known rank-one channel projector.  For the uncompressed
$c$-channel test, replace $\sqrt2$ in \cref{eq:threshold} by $\sqrt{2c}$.
If the known common range is first compressed exactly to the rank-two
covariance range, \cref{eq:threshold} remains valid unchanged.
\end{theorem}

\begin{proof}
The old test's population null holds by \cref{eq:oldmargin}; its coverage
argument bounds rejection by $\alpha$.

For the full test,
\[
 \ip{A}{K}=sM_1-M_2=-\gamma.                                      \tag{5.8}
\]
Suppose the Gaussian event \cref{eq:coverage} holds and a covariance $Q$
inside the test band is full-null feasible.  Put
$\rho=(1+\eta)/(1-\eta)$.  Transitivity of Loewner order gives
\[
 \rho^{-2}K\preceq Q\preceq\rho^2K.                              \tag{5.9}\label{eq:relativeband}
\]
Let $B=K^{1/2}AK^{1/2}$, $p=\tr B_+$, and $m=\tr B_-$.  Then
$p+m=W$ and $m-p=\gamma$.  The support function of
\cref{eq:relativeband} gives
\[
 \ip{A}{Q-K}
 \le(\rho^2-1)p+(1-\rho^{-2})m.                                   \tag{5.10}
\]
Substituting $p=(W-\gamma)/2$ and $m=(W+\gamma)/2$ shows that the right
side is less than $\gamma$ whenever
\[
 \eta<\frac{r}{(\sqrt{1+r}+1)^2}.                                 \tag{5.11}\label{eq:etapower}
\]
Hence $\ip{A}{Q}<0$.  But for a structured
$Q=H_1(\widetilde M)$,
$\ip{A}{Q}=s\widetilde M_1-\widetilde M_2$, which must be
nonnegative under the full null.  Therefore the primal is infeasible on the
coverage event.  Solving \cref{eq:etapower} for $n$ with $d=2$ gives
\cref{eq:threshold}.
\end{proof}

\subsection{Exact rational fixture}

Take
\[
 s=\frac12,\qquad t=\frac34,\qquad w=\frac18.                     \tag{5.12}
\]
Then
\[
 K=\begin{pmatrix}1&3/32\\3/32&9/128\end{pmatrix},\qquad
 \det K=\frac{63}{1024}>0,                                        \tag{5.13}
\]
\[
 s^2M_0-M_2=\frac{23}{128}>0,\qquad
 sM_1-M_2=-\frac3{128}<0,                                        \tag{5.14}
\]
and
\[
 \gamma=\frac3{128},\qquad W^2=\frac{261}{16384},\qquad
 r=\frac3{\sqrt{261}}.                                            \tag{5.15}
\]
These opposing margins and $W^2$ are verified over exact rationals.

\subsection{Full exponent Grams}

For this paragraph define generalized moments
$M_\beta=\int_0^1y^\beta\,d\bm\nu(y)$ for every real $\beta\ge0$.
For any exponents $a_0,\ldots,a_R\ge0$, the support null implies
\[
 \mathcal K_s^A=
 [sM_{a_i+a_j+1}-M_{a_i+a_j+2}]_{i,j=0}^R\succeq0,               \tag{5.16}
\]
because this is the Gram matrix of $y^{a_i}$ in
$L^2(y(s-y)d\nu)$.  Half-integer exponents couple integer and half-time
data.  Separate-grid constraints are principal submatrices and cannot
dominate the full joint Gram.  Equation (5.16) is operationally testable only
when the corresponding exponent blocks have also been acquired; it is not an
additional constraint available from the integer grid $0,\ldots,q$ alone.

\section{Global separation of fixed alternatives}

Let $\cC_{q,c,s}$ be the closed cone of structured covariances $H_q(M)$
compatible with a positive matrix measure on $[0,s]$.  For fixed
$K\notin\cC_{q,c,s}$, set
\[
 \delta(K)=\dist_F(K,\cC_{q,c,s})>0.                              \tag{6.1}
\]

\begin{theorem}[fixed-alternative consistency]\label{thm:fixed}
For every fixed $K\notin\cC_{q,c,s}$, the rejection probability tends to one.
Uniformity holds on sets satisfying $\norm K_{\rm op}\le B$ and
$\dist_F(K,\cC_{q,c,s})\ge\delta_0>0$.  If
\[
 (1-u)^2K\preceq S\preceq(1+u)^2K,\qquad u<1,                    \tag{6.2}
\]
then every $Q$ in the test band lies in a Frobenius ball about $K$ of radius
\[
 \sqrt d\,\Gamma(u,\eta)\norm{K}_{\rm op},                        \tag{6.3}
\]
where
\[
 \Gamma(u,\eta)=\max\left\{
 \frac{(1+u)^2}{(1-\eta)^2}-1,\,
 1-\frac{(1-u)^2}{(1+\eta)^2}\right\}.                            \tag{6.4}
\]
If this radius is less than $\delta(K)$, rejection is forced.
\end{theorem}

\begin{proof}
Combining the event with $aS\preceq Q\preceq bS$ gives scalar lower and
upper Loewner multiples of $K$.  Hence
$\norm{Q-K}_F\le\sqrt d\,\Gamma\norm{K}_{\rm op}$.  If that ball misses the
closed null cone, no structured sequence is feasible.  Choose $t_n\to\infty$
with $t_n/\sqrt n\to0$ and put $u_n=(\sqrt d+t_n)/\sqrt n$.  The corresponding
Gaussian singular-value event has probability at least $1-2e^{-t_n^2/2}$,
while $u_n,\eta_n\to0$.  Hence the displayed radius is eventually below the
fixed positive distance; the same argument is uniform under the two stated
bounds.
\end{proof}

\section{Tangent geometry}

Fix $K^0=H_q(M^0)\succ0$ in the null; its localizer may be singular.  Put
\[
\mathcal V=(\Sym_c)^{2q+1},\quad
\mathsf H(M)=H_q(M),\quad
\mathsf C(M)=C_{q-1,s}(M).                                       \tag{7.1}
\]
The map $\mathsf H$ is injective because every moment occurs on an
anti-diagonal.  Let $U_0$ span $\ker\mathsf C(M^0)$ and define
\[
\cT_0=\{\mathsf H(D):D\in\mathcal V,\,
U_0^*\mathsf C(D)U_0\succeq0\}.                                  \tag{7.2}\label{eq:tangent}
\]

\begin{theorem}[exact tangent preimage]\label{thm:tangent}
The Bouligand tangent cone of $\cC_{q,c,s}$ at $K^0$ equals $\cT_0$.
\end{theorem}

\begin{proof}
For $C^0\succeq0$,
\[
T_{\mathbb S_+}(C^0)=\{E:U_0^*EU_0\succeq0\}.                    \tag{7.3}
\]
Necessity follows by differentiating a feasible sequence and compressing to
$\ker C^0$.  Conversely, choose an interior moment sequence
$M^{\rm int}$ generated by a positive density on $(0,s)$ times $I_c$.
Both $\mathsf H(M^{\rm int})$ and $\mathsf C(M^{\rm int})$ are positive
definite.  Therefore $D^{\rm sl}=M^{\rm int}-M^0$ satisfies
\[
U_0^*\mathsf C(D^{\rm sl})U_0\succ0.                              \tag{7.4}
\]
Robinson's constraint qualification gives
\[
T_{\{M:\mathsf C(M)\succeq0\}}(M^0)
=\{D:U_0^*\mathsf C(D)U_0\succeq0\}.                              \tag{7.5}
\]
The $H_q$ positivity constraint is inactive at $K^0\succ0$.  In fact
$\mathsf C(M^0)+\mathsf C(D^{\rm sl})=\mathsf C(M^{\rm int})\succ0$, so the
affine semidefinite constraint satisfies Robinson's condition.
Applying the injective map $\mathsf H$ proves the claim.
\end{proof}

If $K^0$ is singular, its PSD constraint is active and the experiment lives
on a smaller range.  One must add its kernel compression and exact facial
reduction.  The regular likelihood theorem below is restricted to
$K^0\succ0$.

\section{Wishart LAN and local asymptotic power}

For real zero-mean Gaussian observations, the Fisher metric is
\[
\ip{H}{J}_{I_0}
=\frac12\tr[(K^0)^{-1}H(K^0)^{-1}J].                              \tag{8.1}
\]
Let $I_0^{1/2}$ be a whitening isometry and put
$\mathcal K_0=I_0^{1/2}\cT_0$.

\begin{theorem}[LAN and cone-projection likelihood]\label{thm:lan}
Under
\[
K_n=K^0+\frac{H}{\sqrt n}+o(n^{-1/2}),                             \tag{8.2}
\]
\[
\log\frac{dP_{K_n}^{\otimes n}}{dP_{K^0}^{\otimes n}}
=\ip{Z_n}{H}_{I_0}-\frac12\norm{H}_{I_0}^2+o_{P_{K^0}}(1),         \tag{8.3}
\]
where $Z_n$ is the Fisher--Riesz representation of the central score: for
every fixed $J$, $\ip{Z_n}{J}_{I_0}$ is the normalized score in direction
$J$.  Thus $I_0^{1/2}Z_n\Rightarrow G$, standard Gaussian in $\Sym_d$.
Twice the log-likelihood ratio between the unrestricted covariance model and
the structured null obeys
\[
\boxed{\Lambda_n\Rightarrow
\dist^2(G+I_0^{1/2}H,\mathcal K_0).}                               \tag{8.4}\label{eq:conelimit}
\]
\end{theorem}

\begin{proof}
Up to constants,
\[
\ell_n(K)=-\frac n2\{\log\det K+\tr(K^{-1}S)\}.                    \tag{8.5}
\]
Insert $K=K^0+H/\sqrt n$ and use
\[
\begin{aligned}
\log\det(K^0+E)&=\log\det K^0+\tr((K^0)^{-1}E)
-\tfrac12\tr[((K^0)^{-1}E)^2]+o(\norm E^2),\\
(K^0+E)^{-1}&=(K^0)^{-1}-(K^0)^{-1}E(K^0)^{-1}\\
&\quad +(K^0)^{-1}E(K^0)^{-1}E(K^0)^{-1}+o(\norm E^2).
\end{aligned}                                                       \tag{8.6}
\]
The covariance CLT supplies the Gaussian score and Fisher quadratic term.
The local unrestricted likelihood is
$h\mapsto\ip{G+I_0^{1/2}H}{h}-\norm h^2/2$.  By
\cref{thm:tangent} and linearized Slater, local null sets converge to
$\mathcal K_0$.  Completing the square after unrestricted and constrained
maximization gives \cref{eq:conelimit}.
\end{proof}

For proper complex Gaussian observations the factor $1/2$ in the Fisher
metric is absent, the covariance is complex Wishart, and $G$ is standard in
the real Hilbert space $\Herm_d$; all tangent and polar cones are then taken
in $\Herm_d$.  This is a separate extension, not a change of one scalar factor
in the real experiment.

Let $c_\alpha(\mathcal K_0)$ be an upper $\alpha$ quantile of
$\dist^2(G,\mathcal K_0)$.

\begin{corollary}[local size and power]\label{cor:local}
The cone-likelihood test has asymptotic size at most $\alpha$ for every local
null sequence.  Its local power against $H$ is
\[
\pi_{K^0}(H)=
\Pr\{\dist^2(G+I_0^{1/2}H,\mathcal K_0)>c_\alpha(\mathcal K_0)\}.  \tag{8.7}\label{eq:localpower}
\]
\end{corollary}

\begin{proof}
Every local null limit $h$ belongs to $\mathcal K_0$.  Since the cone is
closed under addition,
\[
\dist(G+h,\mathcal K_0)\le\dist(G,\mathcal K_0)                   \tag{8.8}
\]
pointwise.  The vertex is least favorable.  Randomization at a quantile atom
gives exact asymptotic size; otherwise the test is conservative.
\end{proof}

At $H=0$, Moreau decomposition gives
$\norm{\Pi_{\mathcal K_0^\circ}G}^2$.  For general spectrahedral preimages,
we retain the description as a Gaussian cone-projection law.

\begin{theorem}[matching $n^{-1/2}$ boundary]\label{thm:minimax}
At each regular boundary point, for a fixed outward direction
$V\notin\cT_0$ and a sequence $K_n=K^0+\varepsilon_nV\succ0$ with
$\varepsilon_n=o(n^{-1/2})$, no level-$\alpha$ test has power above
$\alpha+o(1)$.  Displacement
$H/\sqrt n$ has the nontrivial power in \cref{cor:local}; and alternatives
that remain in a compact positive-definite neighborhood and satisfy
$\sqrt n\,\dist_{I_0}(K_n,\cC_{q,c,s})\to\infty$ are detected by the
structured Loewner-feasibility test with probability tending to one.
\end{theorem}

\begin{proof}
For an outward $V\notin\cT_0$,
\[
D(P_{K^0+\varepsilon V}^{\otimes n}\|P_{K^0}^{\otimes n})
=\frac{n\varepsilon^2}{2}\norm{V}_{I_0}^2+o(n\varepsilon^2).       \tag{8.9}
\]
Pinsker proves the lower bound, \cref{thm:lan} gives the local limit, and
\cref{thm:fixed} plus covariance concentration gives the upper bound.
\end{proof}

If $0<\alpha<1/2$, then at a full-dimensional smooth halfspace face
$\mathcal K_0=\{h:\ip{a}{h}\ge0\}$, $\norm a=1$, with outward
$\rho=-\ip{a}{I_0^{1/2}H}>0$, the exact local power is
\[
\boxed{\pi(\rho)=1-\Phi(z_{1-\alpha}-\rho).}                       \tag{8.10}
\]
This is the Neyman--Pearson envelope in the limiting one-dimensional normal
experiment.  At the atom, including $\alpha\ge1/2$, the critical rule must be
randomized.  If the Hankel image is a proper subspace of $\Sym_d$, the formula
applies conditionally inside that structured submodel; the unrestricted test
also contains the orthogonal equality-constraint component.

\section{Auditable semidefinite dual certificates}\label{sec:dual}

The primal test is useful only if rejection can be checked independently of
the floating-point solver that found it.  We therefore spell out its ordinary
semidefinite Farkas dual.  For block matrices $Y=[Y_{ij}]_{i,j=0}^q$ and
$Z=[Z_{ij}]_{i,j=0}^{q-1}$ define
\[
 \begin{aligned}
 [\mathsf H_q^*Y]_k&=\sum_{i+j=k}Y_{ij},\\
 [\mathsf C_{q-1,s}^*Z]_k
 &=s\sum_{i+j+1=k}Z_{ij}-\sum_{i+j+2=k}Z_{ij},
 \end{aligned}                                                     \tag{9.1}
\]
where empty sums are zero and symmetric symmetrization is understood.  These
are the exact adjoints because
$\ip{Y}{H_q(M)}=\sum_k\ip{[\mathsf H_q^*Y]_k}{M_k}$, and similarly
for $\mathsf C$.

\begin{proposition}[checkable rejection certificate]\label{prop:dual}
If there are positive semidefinite block matrices $Y_A,Y_B,Y_C$ satisfying
\[
 \mathsf H_q^*(Y_A-Y_B)+\mathsf C_{q-1,s}^*Y_C=0                 \tag{9.2}\label{eq:dualstationarity}
\]
and
\[
 -a\ip{Y_A}{S}+b\ip{Y_B}{S}<0,                                  \tag{9.3}\label{eq:dualobjective}
\]
then the confidence feasibility problem is infeasible.  Conversely, every
infeasible instance of this bounded structured-band problem admits a
certificate of this form.
\end{proposition}

\begin{proof}
For any primal-feasible $M$, positivity gives
\[
 0\le \ip{Y_A}{H_q(M)-aS}+\ip{Y_B}{bS-H_q(M)}
       +\ip{Y_C}{C_{q-1,s}(M)}.                                  \tag{9.4}
\]
Stationarity cancels every moment block, leaving exactly the strictly
negative number in \cref{eq:dualobjective}, a contradiction.  For the
converse, suppose approximate feasibility had residual tending to zero.  The
sandwich constraints bound $\mathsf H_q(M)$; injectivity of $\mathsf H_q$ then
bounds the moment variables.  A convergent subsequence would be exactly
feasible because the PSD cones are closed, a contradiction.  Thus the affine
image and product PSD cone have positive distance, and strong separation
yields the displayed ordinary multipliers.  This compact sandwich argument
is special to the present bounded band; generic singular semidefinite systems
may require facial reduction or Ramana's extended dual \cite{Ramana1997}.
\end{proof}

\subsection{An exact rational dual certificate}

Collapse the band to $a=b=1$ and set $S=K$ from \cref{eq:twoatom} with the
rational values in \cref{sec:strict}.  Let
\[
 D=\begin{pmatrix}0&-1/4\\-1/4&1\end{pmatrix},\quad
 Y_A=D+I=\begin{pmatrix}1&-1/4\\-1/4&2\end{pmatrix},\quad
 Y_B=I,\quad Y_C=1.                                               \tag{9.5}
\]
Both matrix multipliers are positive definite.  Since $s=1/2$,
the three moment coefficients of
$\mathsf H_1^*D+\mathsf C_{0,s}^*1$ are
\[
 D_{00}=0,\qquad 2D_{01}+s=0,\qquad D_{11}-1=0.                  \tag{9.6}
\]
Thus stationarity holds over the rationals, while
\[
 -\ip{Y_A}{K}+\ip{Y_B}{K}=-\ip D K=-\frac3{128}<0.              \tag{9.7}
\]
This is a solver-independent certificate that the full constraint rejects
the population covariance.  Its entries, principal minors, stationarity
residuals, and objective are replayed exactly by the companion script.

\section{Algorithm and reproducible replay}\label{sec:algorithm}

For exponent depth $q$, channel size $c$, support cutoff $s$, and level
$\alpha$, the operational procedure is:
\begin{enumerate}[label=\arabic*.]
 \item Acquire $n$ independent joint sketch vectors
 $z_j=(z_{j,0},\ldots,z_{j,q})$ and form
 $S=n^{-1}\sum_jz_jz_j^\mathsf T$.
 \item Set $d=(q+1)c$, compute $\eta,a,b$ from
 \cref{eq:coverage}, and stop as sample-size unresolved if $\eta\ge1$.
 \item Introduce real-symmetric variables $M_0,\ldots,M_{2q}$ and solve
 the structured Loewner feasibility problem in Section~4.  Reject if and
 only if infeasible.
 \item On rejection, request PSD multipliers from the solver, rationalize or
 outward-round them, and verify PSD, stationarity, and a strictly negative
 objective independently.  If a solver fails to recover the guaranteed
 ordinary certificate numerically, use facial reduction to expose and
 preserve the relevant face before reconstructing it.
\end{enumerate}

The data path uses a single sample covariance, hence cross-exponent blocks
and their sampling correlations are never discarded.  Complexity is that of
an SDP with $2q+1$ real-symmetric $c\times c$ variables, two PSD constraints of
order $(q+1)c$, and one of order $qc$.  Anti-diagonal consistency is exact
because $H_q(M)$ is assembled from shared variables, rather than imposed by
noisy equality tests.

The lightweight replay file
\texttt{verify\_matricial\_hausdorff.py} uses exact rational arithmetic for
the algebraic fixtures and floating-point normal quantiles only for the
reported Gaussian illustration.  It checks
the even and odd block criteria on noncommuting positive $2\times2$ atomic
weights, vector-polynomial identities, Loewner inversion, an intentionally
corrupted moment witness, the strict two-atom fixture, the dual certificate,
and a regular smooth-face tangent fixture.  It does not simulate coverage, optimize
over scientific models, or certify novelty.  A successful run therefore
certifies the paper's algebraic fixtures, not its empirical scope.

\section{Independent-moment observations as a secondary variant}

Some experiments produce independent covariance estimates
$\widehat M_k$ rather than one joint exponent vector.  If simultaneous bands
\[
 \underline M_k\preceq M_k\preceq\overline M_k,\qquad 0\le k\le m, \tag{11.1}
\]
have joint coverage at least $1-\alpha$, intersecting those bands with the
parity-appropriate LMIs in \cref{thm:cone} gives the same finite-sample level
argument.  Bonferroni, a joint bootstrap with proved coverage, or exact
Wishart bounds may construct the simultaneous event.  This variant is useful
but strictly less informative for the present separation question: estimating
blocks independently obscures the joint covariance geometry and may fail to
retain the half-time cross blocks that make \cref{thm:strict} possible.

The joint model is therefore the primary theorem, not an example of the
independent oracle.  Conversely, joint acquisition does not manufacture
unobserved moments: with exponents $0,\ldots,q$ it exposes exactly sums
$0,\ldots,2q$.  Claims about a larger truncated cone require additional
exponents or assumptions.

\section{Relation to prior work and scope boundaries}\label{sec:scope}

The matricial Hausdorff equivalences themselves are classical
\cite{DetteStudden2002,ChoqueRiveroEtAl2006,ChoqueRiveroEtAl2007,
FritzscheKirsteinMaedler2019}; scalar moment and localizing matrices belong to
the established truncated-moment framework
\cite{Lasserre2001,HeltonNie2012}.  Gaussian covariance concentration,
constrained likelihood limits, and exact rational SDP certification likewise
have substantial precedents
\cite{DavidsonSzarek2001,KoltchinskiiLounici2017,Chernoff1954,Shapiro1985,Geyer1994,
PeyrlParrilo2008}.  We do not rename any of those components as new.

This paper explicitly extends the finite-sample spectral support framework
archived as \texttt{ARR-2026-77QM18J2KG9679B7}
\cite{ARRPrior2026}.  Its increment is the integrated theorem chain:
one operational exponent sketch produces the entire block-Hankel covariance;
the parity-complete matrix cone uses every cross block; a closed rational
family proves strict same-data power over the preceding integer-time
localizer; a single Loewner event gives a multiplicity-free conic test; and
the tangent preimage identifies its sharp $n^{-1/2}$ boundary law.  The
strictness claim is relative to the precisely stated predecessor constraint,
not to all conceivable support tests.

Several limitations are structural.
\begin{itemize}
 \item Gaussian sketches give the stated finite-sample band.  Sub-Gaussian or
 dependent data need a separately proved covariance region.
 \item The cutoff hypothesis concerns existence of a positive matrix measure
 on $[0,s]$.  It does not identify a unique measure, dynamical generator, or
 physical microscopic mechanism.
 \item Finite truncation cannot detect every alternative distribution.  The
 consistency theorem is against covariances outside the exposed truncated
 cone, not against alternatives sharing all observed blocks with a null law.
 \item The LAN theorem assumes a positive-definite covariance and the stated
 linearized Slater condition.  More singular intersections can have different
 rates and require stratified facial analysis.
 \item Ordinary floating-point infeasibility is not a proof.  A preserved dual
 certificate, preferably rational or interval verified, is required for an
 auditable individual rejection.
\end{itemize}

\section{Conclusion}

A joint half-time exponent sketch turns a matrix-valued support question into
a single structured Wishart problem.  The full block-Hankel covariance and
its Hausdorff localizer yield finite-sample validity, strict same-data power,
global fixed-alternative consistency, and an exact Gaussian tangent-cone
limit.  The rational fixture isolates why the gain occurs: the discarded
cross moment is negative in the support localizer even while the older
integer-grid condition is strictly feasible.  The result is thus stronger
than an additional moment inequality: it gives an operational acquisition,
a sharp boundary theory, and independently checkable rejection witnesses.

\printbibliography
\end{document}
